\begin{thebibliography}{99}
\addcontentsline{toc}{chapter}{参考文献}
\small\raggedright\setlength{\itemsep}{5pt}\setlength{\parskip}{0pt}
\bibitem{haus} H. A. Haus and J. R. Melcher. \emph{Electromagnetic Fields and Energy}. Prentice Hall, 1989. MIT 开放教材，涵盖准静态、磁化、磁扩散、能量与波动。 \href{https://web.mit.edu/6.013_book/www/book.html}{MIT 在线教材}.

\bibitem{itisdielectric} IT'IS Foundation. \emph{Tissue Properties Database}. 随频率变化的介电参数及原始文献。 \href{https://itis.swiss/virtual-population/tissue-properties/database}{官方数据库}.

\bibitem{itislf} IT'IS Foundation. \emph{Low Frequency (Conductivity)}. 组织电导率、统计分布与数据来源。 \href{https://itis.swiss/virtual-population/tissue-properties/database/low-frequency-conductivity}{官方低频数据汇编}.

\bibitem{simnibs} SimNIBS developers. \emph{SimNIBS Documentation}. 经颅电刺激与磁刺激的有限元仿真。 \href{https://simnibs.github.io/simnibs/}{官方文档}. 具体方法及版本对应文献，请参阅文档中的论文列表。

\bibitem{cem} E. Somersalo, M. Cheney, and D. Isaacson. ``Existence and Uniqueness for Electrode Models for Electric Current Computed Tomography.'' \emph{SIAM Journal on Applied Mathematics}, 52(4):1023--1040, 1992. \href{https://doi.org/10.1137/0152060}{doi:10.1137/0152060}.

\bibitem{comsolbh} L. Liu. ``How the B-H Curve Affects a Magnetic Analysis (and How to Improve It).'' COMSOL Blog, 26 November 2019. \href{https://www.comsol.com/blogs/how-the-b-h-curve-affects-a-magnetic-analysis-and-how-to-improve-it}{官方文章}. 本讲义中的矢量微分张量由\chref{ch:femvector}的显式局部线性化得到。

\bibitem{femm} D. Meeker. \emph{Finite Element Method Magnetics: User's Manual}. 材料输入、B--H 插值、永磁体与力的计算。 \href{https://www.femm.info/wiki/Documentation/}{FEMM 文档}.

\bibitem{mfem} MFEM developers. \emph{MFEM Examples}. 包含 Poisson 与 Maxwell 类问题、N\'ed\'elec 空间以及混合表述。 \href{https://mfem.org/examples/}{官方示例}.

\bibitem{getdp} P. Dular and C. Geuzaine. \emph{GetDP Reference Manual}. 方程表述、函数空间、约束与电磁学示例。 \href{https://getdp.info/doc/texinfo/getdp.html}{官方手册}.

\bibitem{rf} C. M. Collins and Z. Wang. ``Calculation of radiofrequency electromagnetic fields and their effects in MRI of human subjects.'' \emph{Magnetic Resonance in Medicine}, 65(5):1470--1482, 2011. \href{https://doi.org/10.1002/mrm.22845}{doi:10.1002/mrm.22845}.

\bibitem{concomitant} M. A. Bernstein, X. J. Zhou, J. A. Polzin, K. F. King, A. Ganin, N. J. Pelc, and G. H. Glover. ``Concomitant gradient terms in phase contrast MR: Analysis and correction.'' \emph{Magnetic Resonance in Medicine}, 39(2):300--308, 1998. \href{https://doi.org/10.1002/mrm.1910390218}{doi:10.1002/mrm.1910390218}.

\bibitem{bfieldtheory} A. J. M\"akinen, R. Zetter, J. Iivanainen, K. C. J. Zevenhoven, L. Parkkonen, and R. J. Ilmoniemi. ``Magnetic-field modeling with surface currents. Part I. Physical and computational principles of bfieldtools.'' \emph{Journal of Applied Physics}, 128:063906, 2020. \href{https://doi.org/10.1063/5.0016090}{doi:10.1063/5.0016090}.

\bibitem{bfieldsoftware} R. Zetter, A. J. M\"akinen, J. Iivanainen, K. C. J. Zevenhoven, R. J. Ilmoniemi, and L. Parkkonen. ``Magnetic field modeling with surface currents. Part II. Implementation and usage of bfieldtools.'' \emph{Journal of Applied Physics}, 128:063905, 2020. \href{https://doi.org/10.1063/5.0016087}{doi:10.1063/5.0016087}.

\bibitem{rattay} F. Rattay. ``Analysis of models for external stimulation of axons.'' \emph{IEEE Transactions on Biomedical Engineering}, BME-33(10):974--977, 1986.

\bibitem{davids} M. Davids, B. Gu\'erin, A. vom Endt, L. R. Schad, and L. L. Wald. ``Prediction of peripheral nerve stimulation thresholds of MRI gradient coils using coupled electromagnetic and neurodynamic simulations.'' \emph{Magnetic Resonance in Medicine}, 81(1):686--701, 2019. \href{https://doi.org/10.1002/mrm.27382}{doi:10.1002/mrm.27382}.

\bibitem{polarization} H. A. Haus and J. R. Melcher. ``Polarization Density,'' Section 6.1 of \emph{Electromagnetic Fields and Energy}. \href{https://web.mit.edu/6.013_book/www/chapter6/6.1.html}{MIT 教材：电偶极子与束缚电荷}.

\bibitem{magnetization} H. A. Haus and J. R. Melcher. ``Magnetization Density,'' Section 9.1 of \emph{Electromagnetic Fields and Energy}. \href{https://web.mit.edu/6.013_book/www/chapter9/9.1.html}{MIT 教材：磁矩与磁化}.

\bibitem{vectorpoisson} H. A. Haus and J. R. Melcher. ``The Vector Potential and the Vector Poisson Equation,'' Section 8.1. \href{https://web.mit.edu/6.013_book/www/chapter8/8.1.html}{MIT 教材：规范、源积分与二维势函数}.

\bibitem{vectorboundary} H. A. Haus and J. R. Melcher. ``Vector Potential and the Boundary Value Point of View,'' Section 8.6. \href{https://web.mit.edu/6.013_book/www/chapter8/8.6.html}{MIT 教材：磁矢势环量与磁通}.

\bibitem{demagnetizing} R. Fitzpatrick. ``Uniformly Magnetized Sphere,'' in \emph{Classical Electromagnetism}. 得克萨斯大学奥斯汀分校讲义。 \href{https://farside.ph.utexas.edu/teaching/jk1/Electromagnetism/node61.html}{内部退磁场的推导}.

\bibitem{qslimits} H. A. Haus and J. R. Melcher. ``Conditions for Fields to be Quasistatic,'' Section 3.3. \href{https://web.mit.edu/6.013_book/www/chapter3/3.3.html}{MIT 教材：保留的场与忽略的反馈}.

\bibitem{mfemem} MFEM developers. ``Electromagnetics Mini Applications.'' \href{https://mfem.org/electromagnetics/}{官方场方程与有限元示例}.

\bibitem{simnibsfields} SimNIBS developers. \texttt{simnibs.simulation.fem.calc\_fields}, 官方 API 文档。 \href{https://simnibs.github.io/simnibs/build/html/documentation/python_lib/auto/simnibs.simulation.fem.calc_fields.html}{场重建、单位及输出符号约定}.


\bibitem{griffiths} D. J. Griffiths. \emph{Introduction to Electrodynamics}, 4th ed. Cambridge University Press, 2017. 第 4、6 章：极化、磁化、束缚电荷与束缚电流、$\vb D$ 与 $\vb H$。

\bibitem{purcell} E. M. Purcell and D. J. Morin. \emph{Electricity and Magnetism}, 3rd ed. Cambridge University Press, 2013. 第 10、11 章：物质中的电场与磁场，偶极子内部的场。


\bibitem{fitzgerald} A. E. Fitzgerald, C. Kingsley Jr., and S. D. Umans. \emph{Electric Machinery}, 7th ed. McGraw-Hill, 2014. 第 1、3 章：磁路、磁链与电感、磁能与共能、虚功法求力、永磁体与负载线。

\bibitem{osborn} J. A. Osborn. ``Demagnetizing Factors of the General Ellipsoid.'' \emph{Physical Review}, 67(11--12):351--357, 1945. \href{https://doi.org/10.1103/PhysRev.67.351}{doi:10.1103/PhysRev.67.351}.

\bibitem{fritsch} F. N. Fritsch and R. E. Carlson. ``Monotone Piecewise Cubic Interpolation.'' \emph{SIAM Journal on Numerical Analysis}, 17(2):238--246, 1980. \href{https://doi.org/10.1137/0717021}{doi:10.1137/0717021}.

\bibitem{nedelec} J.-C. Nédélec. ``Mixed Finite Elements in $\mathbb{R}^3$.'' \emph{Numerische Mathematik}, 35(3):315--341, 1980. \href{https://doi.org/10.1007/BF01396415}{doi:10.1007/BF01396415}.

\bibitem{oberkampf} W. L. Oberkampf and C. J. Roy. \emph{Verification and Validation in Scientific Computing}. Cambridge University Press, 2010. ISBN 978-0-521-11360-1.

\end{thebibliography}
